\begin{env}[7.3.4]
\label{IV.2.s7.3.4-fr}
En lo que sigue supondremos siempre que la propiedad $\boldsymbol Q$
es tal que se cumplen las tres condiciones siguientes:
\begin{enumerate}
\item[$(\mathrm P_{\mathrm I})$] \emph{(transitividad).} --- Si
$f:X\to Y$ es un morfismo regular (6.8.1) y $g:Y\to Z$ un
$\boldsymbol P$-morfismo, entonces $g\circ f$ es un
$\boldsymbol P$-morfismo.
\item[$(\mathrm P_{\mathrm{II}})$] \emph{(descenso).} --- Si
$f:X\to Y$ y $g:Y\to Z$ son morfismos de preesquemas localmente
noetherianos tales que $f$ es fielmente plano y $g\circ f$ es un
$\boldsymbol P$-morfismo, entonces $g$ es un
$\boldsymbol P$-morfismo.
\item[$(\mathrm P_{\mathrm{III}})$.] --- Para todo cuerpo $k$, se cumple
la propiedad $\boldsymbol P(\operatorname{Spec}(k),k)$.
\end{enumerate}
\end{env}
